\chap{目标与基本概念：开发系统作为能力放大器}
\label{sec:framework}
\subsection{能力放大的含义}
本文所说的能力放大，是组织方法、工具和积累使同等模型能够完成更复杂或质量更高的工程工作，并使更强模型能够进一步扩展这种范围。它包含同一能力条件下的增益，也包含能力提升后的可扩展性。判断时需要控制项目条件，观察真实产物与人的投入，不能用成员数量、执行轮数或生成代码长度替代效果。

放大的第一种来源是有效组合。不同成员深入相对独立的职责范围，共享明确的协议，在必要的地方协作。第二种来源是跨时间积累：设计理由、故障案例、工具和验证经验降低后续调查成本。第三种来源是现实反馈：交付、测试和使用让判断接触实际结果，减少重复的错误假设。第四种来源是交互质量：人更容易理解关键问题，因而能以较少努力给出有效引导。

这些来源具有条件性。并行会增加协调，记忆会过时，反馈可能不完整，生成界面也可能误导。一个有用的框架必须允许发现负收益并调整结构。某项目只需要一个足够强的 Agent 时，持续身份、证据记录和交互投影仍可使用；团队划分应服从实际工作，而非成为进入框架的先决条件。

\subsection{核心对象与职责}
\begin{table}[htbp]\centering\small
\caption{框架中的核心对象及其持续职责}\label{tab:objects}
\begin{tabularx}{\linewidth}{p{2.4cm}YY}\toprule
对象 & 含义 & 需要保持的联系\\\midrule
开发 Agent & 能理解任务并通过工具采取行动的执行主体 & 动作与目标、结果及执行环境\\
持续成员 & 跨模型会话保有身份与职责的 Agent & 维护范围、经验、未完工作\\
团队 & 围绕共同目标协作的一组成员 & 委派、协作关系与结果责任\\
软件模块 & 围绕相对稳定职责组织的实现与契约 & 内部决定与外部可依赖行为\\
语义规范 & 对目标、约束与成功含义的可理解表达 & 当前意图与实现、验证依据\\
工作证据 & 对特定内容和环境进行观察或检查的结果 & 范围、时间、适用条件与来源\\
交互投影 & 面向当前理解任务组织的视图与操作 & 显示对象、事实来源与操作效果\\\bottomrule
\end{tabularx}\end{table}

在 Portal 的术语中，持续成员称为 Resident，团队及其协作关系由 Lineage 承载，面向人的工作界面称为 Workbench。本文使用这些名称描述具体实例，同时保留一般定义。模型会话是成员使用的一种执行资源；更换会话不应自然地丢失长期职责。工作界面也只是当前任务的一种表达，同一组工作事实可以支持不同视图。

身份、知识和上下文之间需要区分。身份回答谁对某范围持续负责，知识保存经过整理的可复用内容，上下文是某次判断能够直接访问的信息。将所有历史塞进每个会话会增加成本，也会让当前任务被无关细节淹没。组织的作用之一，是让需要的事实在适当时刻进入适当成员的上下文。

\subsection{三个相互连接的工作过程}
框架包含目标解释、工程执行和人机反馈三个过程。目标解释将人的表达联系到当前软件与已有约束，形成可开展的工作。工程执行组织实现、验证和交付。人机反馈选择值得呈现的变化，帮助人理解、体验和调整方向。三个过程可以并行推进，并通过稳定对象和记录协调。

例如，某成员已经在改善取消响应时，人可能反馈“真正关心的是取消后还能不能查看中间结果”。新的目标会改变输出保留策略，却不必使已经完成的状态观察工作全部失效。团队可以识别受影响的契约与验证范围，继续保留仍然成立的成果。持续演化要求这种局部更新能力，而不仅是完整流程的重复执行。

\fig{02-amplification}{能力放大的四个来源与三类成本。各作用之间存在条件关系；收益需通过产物质量、变化适应能力和人的实际投入共同观察。图中关系是研究假设，不表示已测得的因果效应。}{amplification}

\subsection{人的控制成本与系统资源}
为便于讨论，可把人的参与成本表示为
\begin{equation}
 C_H = T_{\text{理解}}+T_{\text{比较}}+T_{\text{表达}}+T_{\text{纠偏}}+T_{\text{恢复}},
\end{equation}
其中各项是相应活动占用的时间，恢复包括中断之后重建任务背景。这个分解用于提示观察项目，不假定不同人的每一分钟具有相同价值，也不要求所有项目都进行精密计时。重要的是发现成本究竟发生在哪里：信息已经展示出来，仍可能需要大量恢复和解释工作。

系统侧成本包含模型推理、工具执行、等待、数据维护及协调。token 是其中一种可记录资源，并不直接等于金钱、能耗或认知工作量。不同模型的分词、缓存、输入输出价格和执行速度不同，比较时应同时报告模型配置、时间与费用口径。将人类时间和机器费用压成单个数字，需要明确权重及适用目的。

质量也需要多个维度。任务结果正确、用户体验顺畅、变化影响可控制、系统可维护，可能并不同步改善。本文主张在给定质量与资源约束下比较开发方式，必要时呈现多目标的可选结果。一个方式节省 token 却显著增加用户理解时间，不能仅凭模型成本更低就判断更优。

\subsection{能力边界与收益发生的条件}
可承担的复杂度包含多个方面：涉及多少独立职责、依赖之间是否存在并发与时间关系、需求变化能否局部处理，以及系统需要连续维护多久。一个文件很多但行为简单的项目，未必比文件较少却包含复杂状态协作的项目更难。比较能力边界时，应具体描述变化结构与质量要求。

组织增益可能随任务规模先上升后下降。少量独立工作可以从并行中受益，过多成员则可能增加协议协调和重复调查。长期经验也只有在新任务与既有职责存在联系时才容易复用。这意味着规模扩大并不自动带来更高放大率，研究需要观察收益的适用区间。

可以把一次开发组织选择理解为在质量约束下寻找合适的资源组合。某任务最需要减少人的理解时间，另一任务则受运行工具延迟限制。系统先识别主要限制，再调整成员、验证和表达投入。统一框架提供可选机制，具体选择由实际瓶颈决定。

\subsection{组织自主性与可理解性}
自主性意味着成员在目标和授权范围内选择行动、分工与检查方式。可理解性要求人在需要时能够获得足以判断的事实和解释。两者可以相互支持：清楚的目标与低成本反馈减少日常微观指挥，真实可查的结果使人能够接受更大的委派范围。

这里需要提供的是工作层面的解释，包括依据哪些约束、比较过哪些可执行方案、采取了什么操作以及观察到什么结果。系统无需把模型内部所有生成过程作为汇报材料。有效解释应能帮助人作出下一步判断，并与可核查的工作对象联系。

当模型能力增强时，一些原先需要多个成员协作的任务可能被一个成员承担。与此同时，项目也可能扩展到更长维护周期和更多依赖。框架的稳定部分是目标、证据、反馈和连续性；成员数量、上下文策略和表达形式则可以变化。持续放大能力意味着反复检验这种组合是否仍然带来工程收益。
